\chapter{Crittografia a Chiave Pubblica e Autenticazione del Messaggio}

\section{Funzioni di Hash Sicure}

La funzione di hash unidirezionale, o \textit{funzione di hash sicura}, è un elemento crittografico fondamentale non solo per l'autenticazione dei messaggi, ma anche per le firme digitali. Essa consente di ottenere un'impronta digitale (\textit{digest}) di lunghezza fissa a partire da un input di lunghezza arbitraria, garantendo efficienza computazionale e compattezza del risultato. 

\begin{defbox}[Proprietà di una Funzione di Hash Sicura]
Le funzioni di hash moderne sono progettate per soddisfare proprietà di sicurezza ben precise:
\begin{itemize}
    \item \textbf{Resistenza alla preimmagine (One-way):} Dato l'hash $h$, è computazionalmente impraticabile risalire al messaggio originale $M$.
    \item \textbf{Resistenza alla seconda preimmagine:} Dato un messaggio $M_1$, è impraticabile trovare un $M_2 \neq M_1$ tale che $H(M_1) = H(M_2)$.
    \item \textbf{Resistenza alle collisioni:} È impraticabile trovare una qualsiasi coppia di messaggi distinti $(M_1, M_2)$ che producano lo stesso valore di hash $H(M_1) = H(M_2)$.
\end{itemize}
\end{defbox}

\subsection{Funzioni di Hash Semplici}
Tutte le funzioni di hash operano secondo principi generali comuni. L'input (messaggio, file, ecc.) viene visto come una sequenza di blocchi di $n$ bit. L'input viene elaborato un blocco alla volta in modo iterativo per produrre un valore di hash di $n$ bit.

Una delle funzioni di hash più semplici è l'operazione di XOR (exclusive-OR) bit a bit di ogni blocco:
$$C_i = b_{i1} \oplus b_{i2} \oplus \dots \oplus b_{im}$$
dove:
\begin{itemize}
    \item $C_i$ è l'$i$-esimo bit del codice hash.
    \item $m$ è il numero di blocchi da $n$ bit.
    \item $b_{ij}$ è l'$i$-esimo bit nel $j$-esimo blocco.
\end{itemize}

Questa operazione (controllo di ridondanza longitudinale) calcola una semplice parità. È ragionevolmente efficace per rilevare errori accidentali, ma inefficace con dati formattati in modo prevedibile.

\begin{warnbox}[Inadeguatezza del RXOR per la Sicurezza]
Una semplice miglioria è la procedura RXOR (rotazione di 1 bit del valore di hash per ogni blocco processato). Sebbene fornisca una buona misura di integrità contro l'alterazione casuale, è \textbf{inutile per la sicurezza crittografica}: un attaccante può calcolare facilmente un blocco finale ad hoc per generare qualsiasi hash desiderato.
\end{warnbox}

\subsection{La Funzione di Hash Sicura SHA}
Negli ultimi anni, la famiglia di funzioni di hash più utilizzata è stata la Secure Hash Algorithm (SHA).
\begin{description}
    \item[SHA-1:] Pubblicata nel 1995, produce un hash di 160 bit. Oggi è deprecata a causa di vulnerabilità note alle collisioni.
    \item[SHA-2:] Pubblicata nel 2002, è una famiglia con diverse lunghezze di output (SHA-224, SHA-256, SHA-384, SHA-512). Sono attualmente considerate altamente sicure e rappresentano lo standard industriale.
    \item[SHA-3:] Pubblicata nel 2015 dal NIST, ha una struttura interna \textit{sponge construction} completamente diversa da SHA-2, offrendo un'alternativa sicura nel caso venissero scoperte vulnerabilità strutturali nella famiglia SHA-2.
\end{description}

\subsection{Elaborazione di SHA-512}
L'algoritmo SHA-512 prende in input un messaggio di lunghezza inferiore a $2^{128}$ bit e produce un digest di 512 bit, elaborando blocchi sequenziali da 1024 bit.

Il processo si articola in 5 fasi:
\begin{enumerate}
    \item \textbf{Padding:} Il messaggio viene riempito affinché la sua lunghezza sia congruente a 896 modulo 1024 (inserendo un bit '1' seguito dai necessari '0').
    \item \textbf{Aggiunta della lunghezza:} Un blocco di 128 bit, indicante la lunghezza originale del messaggio, viene accodato. Il messaggio è ora un multiplo esatto di 1024 bit.
    \item \textbf{Inizializzazione del buffer:} Un buffer di 512 bit (otto registri a 64 bit) viene inizializzato con costanti derivate dalle radici quadrate dei primi 8 numeri primi.
    \item \textbf{Elaborazione in blocchi:} Il cuore è una funzione di compressione a 80 round. Per ogni blocco, la funzione prende il buffer del passo precedente ($H_{i-1}$) e produce il nuovo ($H_i$). L'output dell'80-esimo round viene sommato modulo $2^{64}$ al buffer di input.
    \item \textbf{Output:} Dopo l'elaborazione di tutti i blocchi, il contenuto finale del buffer rappresenta il message digest di 512 bit.
\end{enumerate}

\section{HMAC}

Per usare una funzione di hash per l'autenticazione dei messaggi è necessaria una chiave segreta. L'approccio standardizzato è l'\textbf{HMAC} (Hash-based Message Authentication Code).

\begin{notebox}[Vantaggi dell'HMAC]
Si preferiscono i MAC basati su hash perché le funzioni di hash software sono generalmente più veloci dell'implementazione in software di cifrari simmetrici a blocchi, e le librerie crittografiche sono ampiamente diffuse e ottimizzate.
\end{notebox}

L'algoritmo HMAC combina una chiave segreta con il messaggio tramite due successivi passaggi della funzione di hash:
$$\text{HMAC}(K, M) = H[(K^+ \oplus \text{opad}) \mathbin{\|} H[(K^+ \oplus \text{ipad}) \mathbin{\|} M]]$$

dove:
\begin{itemize}
    \item $H$ è la funzione di hash (es. SHA-256 o SHA-512).
    \item $M$ è il messaggio in chiaro.
    \item $K$ è la chiave segreta condivisa.
    \item $K^+$ è la chiave $K$ allungata con zeri fino alla dimensione del blocco interno della funzione $H$.
    \item \text{ipad} e \text{opad} sono stringhe costanti esadecimali predefinite.
    \item $\|$ rappresenta la concatenazione.
\end{itemize}

La sicurezza di HMAC è formalmente dimostrata ed è strettamente legata alla robustezza della funzione di hash $H$ sottostante. Senza la conoscenza della chiave $K$, non è possibile lanciare attacchi offline né generare MAC validi.

\begin{center}
\begin{tikzpicture}[
    node distance=1.5cm and 1cm,
    box/.style={rectangle, draw=blue!80!black, fill=blue!10, thick, minimum width=2.5cm, minimum height=1cm, align=center, rounded corners},
    xor/.style={circle, draw=black, thick, minimum size=0.6cm, inner sep=0pt},
    arrow/.style={-{Stealth}, thick}
]

\node[box] (k) {\textbf{Chiave $K^+$}};
\node[xor, right=of k] (x1) {$\oplus$};
\node[above=of x1, font=\ttfamily] (ipad) {ipad};
\node[box, right=of x1, fill=yellow!20] (concat1) {Concatena $\parallel$\\Messaggio $M$};
\node[box, right=of concat1, fill=red!20] (h1) {\textbf{Hash $H$}};

\draw[arrow] (k) -- (x1);
\draw[arrow] (ipad) -- (x1);
\draw[arrow] (x1) -- (concat1);
\draw[arrow] (concat1) -- (h1);

\node[xor, below=of x1, yshift=-0.5cm] (x2) {$\oplus$};
\node[below=of x2, font=\ttfamily] (opad) {opad};
\draw[arrow] (k.south) |- (x2.west);
\draw[arrow] (opad) -- (x2);

\node[box, right=of x2, fill=yellow!20] (concat2) {Concatena $\parallel$\\Risultato Hash};
\node[box, right=of concat2, fill=red!20] (h2) {\textbf{Hash $H$}};
\node[box, right=of h2, fill=green!20] (mac) {\textbf{HMAC}};

\draw[arrow] (x2) -- (concat2);
\draw[arrow] (h1.south) -- ++(0,-0.3) -| (concat2.north);
\draw[arrow] (concat2) -- (h2);
\draw[arrow] (h2) -- (mac);

\end{tikzpicture}
\captionof*{figure}{Schema dell'algoritmo HMAC}
\end{center}

\section{Cifratura Autenticata (AE)}

La cifratura autenticata (\textit{Authenticated Encryption}) è una tecnica che garantisce contemporaneamente sia la confidenzialità (cifratura) sia l'integrità/autenticità (MAC) dei dati all'interno di un unico schema ottimizzato.

Un esempio efficiente di tale tecnica è la modalità \textbf{Offset Codebook (OCB)}. OCB trasforma un cifrario a blocchi (come AES) in uno schema AE:
$$C[i] = E_K(M[i] \oplus Z[i]) \oplus Z[i]$$

dove $E_K$ è la cifratura con chiave $K$, e $Z[i]$ è un offset unico e dipendente dalla posizione generato da un \textit{nonce}. L'uso degli offset garantisce che blocchi plaintext identici generino ciphertext diversi, nascondendo i pattern. OCB produce contestualmente un tag di autenticazione complessivo cifrando un checksum crittografico di tutti i blocchi plaintext.

\section{L'Algoritmo di Cifratura a Chiave Pubblica RSA}

Sviluppato nel 1977 da Rivest, Shamir e Adleman, l'RSA è il cifrario asimmetrico a chiave pubblica più diffuso al mondo. È fondamentalmente un cifrario a blocchi in cui sia il plaintext che il ciphertext sono numeri interi compresi tra $0$ e $n-1$.

\subsection{Descrizione dell'Algoritmo}
Le operazioni matematiche di cifratura ($C$) e decifratura ($M$) sono:
$$C = M^e \pmod{n}$$
$$M = C^d \pmod{n} = (M^e)^d \pmod{n} = M^{ed} \pmod{n}$$

Le chiavi sono definite come:
\begin{itemize}
    \item \textbf{Chiave Pubblica (pubblica a tutti):} $PU = \{e, n\}$
    \item \textbf{Chiave Privata (segreta):} $PR = \{d, n\}$
\end{itemize}

\subsection{Generazione delle Chiavi}
\begin{enumerate}
    \item Scegliere due numeri primi molto grandi e distinti, $p$ e $q$.
    \item Calcolare $n = p \times q$.
    \item Calcolare la funzione di Eulero $\phi(n) = (p-1)(q-1)$.
    \item Scegliere l'esponente pubblico $e$ tale che $1 < e < \phi(n)$ e $\gcd(\phi(n), e) = 1$.
    \item Calcolare l'esponente privato $d$ tale che $d \equiv e^{-1} \pmod{\phi(n)}$, ovvero $e \cdot d \equiv 1 \pmod{\phi(n)}$.
\end{enumerate}

\begin{exbox}[Esempio di Calcolo RSA]
Siano scelti i piccoli numeri primi $p=17$ e $q=11$:
\begin{itemize}
    \item Si calcola $n = 17 \times 11 = 187$ e $\phi(n) = 16 \times 10 = 160$.
    \item Si sceglie $e=7$ (coprimo con 160), e si calcola $d=23$ (poiché $23 \times 7 = 161 \equiv 1 \pmod{160}$).
\end{itemize}
Per cifrare il messaggio $M=88$: $C = 88^7 \pmod{187} = 11$.\\
Per decifrare il testo $C=11$: $M = 11^{23} \pmod{187} = 88$.
\end{exbox}

\subsection{La Sicurezza di RSA}
Gli attacchi contro RSA si dividono in varie categorie:
\begin{itemize}
    \item \textbf{Forza Bruta:} Provare tutte le chiavi private. Mitigato usando chiavi molto lunghe.
    \item \textbf{Attacchi Matematici (Fattorizzazione):} Si basano sulla difficoltà computazionale di fattorizzare $n$ nei suoi fattori $p$ e $q$, essenziale per calcolare $\phi(n)$ e ricavare $d$. Attualmente si raccomandano chiavi $n$ comprese tra i 2048 e i 4096 bit.
    \item \textbf{Timing Attacks:} Permettono all'attaccante di dedurre frammenti della chiave misurando con precisione il tempo impiegato per l'esponenziazione modulare. Contromisure includono \textit{blinding} (mascherare il ciphertext prima della computazione).
\end{itemize}

\section{Diffie-Hellman e Altri Algoritmi Asimmetrici}

Il protocollo Diffie-Hellman è stato storicamente il primo algoritmo a chiave pubblica ed è dedicato specificamente allo \textbf{scambio sicuro delle chiavi} per la crittografia simmetrica su canali insicuri, e non alla cifratura di messaggi generici.

\begin{notebox}[Logaritmo Discreto]
La sicurezza di Diffie-Hellman si basa sulla difficoltà computazionale del calcolo dei logaritmi discreti in un campo finito: trovare l'esponente $i$ tale che $b \equiv a^i \pmod{p}$ risulta infattibile per grandi valori.
\end{notebox}

\subsection{L'Algoritmo Diffie-Hellman}
Siano $q$ un numero primo grande e $\alpha$ una sua radice primitiva, parametri pubblici noti ad entrambi.
\begin{enumerate}
    \item L'utente Alice sceglie un numero casuale segreto $X_A < q$ e calcola $Y_A = \alpha^{X_A} \pmod{q}$.
    \item L'utente Bob sceglie un numero casuale segreto $X_B < q$ e calcola $Y_B = \alpha^{X_B} \pmod{q}$.
    \item Alice e Bob si scambiano pubblicamente i risultati $Y_A$ e $Y_B$ sulla rete.
    \item Alice calcola la chiave condivisa: $K = (Y_B)^{X_A} \pmod{q}$.
    \item Bob calcola la chiave condivisa: $K = (Y_A)^{X_B} \pmod{q}$.
\end{enumerate}

Entrambi ottengono la medesima chiave segreta condivisa $K$ poiché:
$$(Y_B)^{X_A} = (\alpha^{X_B})^{X_A} = \alpha^{X_B X_A} = (\alpha^{X_A})^{X_B} = (Y_A)^{X_B} \pmod{q}$$

\begin{center}
\begin{tikzpicture}[
    node distance=1cm and 4cm,
    box/.style={rectangle, draw=blue!80!black, fill=blue!5, thick, minimum width=4cm, minimum height=1.5cm, align=center, rounded corners},
    arrow/.style={-{Stealth}, thick, shorten >=2pt, shorten <=2pt}
]

\node[box, fill=purple!10] (A_sec) {\textbf{Alice} (Conosce $q, \alpha$)\\Sceglie Segreto $X_A$\\$Y_A = \alpha^{X_A} \pmod{q}$};
\node[box, fill=purple!10, right=of A_sec] (B_sec) {\textbf{Bob} (Conosce $q, \alpha$)\\Sceglie Segreto $X_B$\\$Y_B = \alpha^{X_B} \pmod{q}$};

\draw[arrow] (A_sec.east) ++(0, 0.2) -- node[above, font=\small] {Invia Valore Pubblico $Y_A$} ($(B_sec.west)+(0,0.2)$);
\draw[arrow] (B_sec.west) ++(0, -0.2) -- node[below, font=\small] {Invia Valore Pubblico $Y_B$} ($(A_sec.east)+(0,-0.2)$);

\node[box, below=of A_sec, fill=green!20, yshift=0.5cm] (A_key) {\textbf{Generazione Chiave}\\$K = (Y_B)^{X_A} \pmod{q}$};
\node[box, below=of B_sec, fill=green!20, yshift=0.5cm] (B_key) {\textbf{Generazione Chiave}\\$K = (Y_A)^{X_B} \pmod{q}$};

\draw[arrow] (A_sec) -- (A_key);
\draw[arrow] (B_sec) -- (B_key);

\end{tikzpicture}
\captionof*{figure}{Scambio di chiavi Diffie-Hellman}
\end{center}
